\chapter{Principal, Gaussian and Mean Curvature}
\label{ch:vii17-principal-gaussian-mean-curvature}

The shape operator \(S_p\) is a self-adjoint endomorphism of the tangent plane. Its eigenvalues contain the fundamental local bending data:
\[
\boxed{k_1,k_2,\qquad K=k_1k_2=\det S,\qquad
H=\frac{k_1+k_2}{2}=\frac12\operatorname{tr}S.}
\]

\section{Principal curvatures and directions}

\begin{definition}[Principal curvatures]\label{def:vii17-principal}
The eigenvalues \(k_1,k_2\) of \(S_p\) are the principal curvatures. Corresponding unit eigenvectors are principal directions.
\end{definition}

\begin{theorem}[Euler formula]\label{thm:vii17-euler}
If \(e_1,e_2\) is an orthonormal principal basis and
\[
v=\cos\theta\,e_1+\sin\theta\,e_2,
\]
then the normal curvature is
\[
k_n(v)=k_1\cos^2\theta+k_2\sin^2\theta.
\]
\end{theorem}

\begin{proof}
Use \(Sv=k_1\cos\theta\,e_1+k_2\sin\theta\,e_2\) and take the inner product with \(v\).
\end{proof}

\begin{corollary}[Extremal property]\label{cor:vii17-extremal}
The principal curvatures are the minimum and maximum values of normal curvature among unit tangent directions.
\end{corollary}

\section{Gaussian and mean curvature}

\begin{definition}[Gaussian curvature]\label{def:vii17-gaussian}
\[
K(p)=\det S_p=k_1k_2.
\]
\end{definition}

\begin{definition}[Mean curvature]\label{def:vii17-mean}
\[
H(p)=\frac12\operatorname{tr}S_p=\frac{k_1+k_2}{2}.
\]
\end{definition}

\begin{proposition}[Orientation dependence]\label{prop:vii17-orientation}
Reversing the orientation sends
\[
k_i\mapsto-k_i,\qquad H\mapsto-H,\qquad K\mapsto K.
\]
Thus Gaussian curvature is independent of orientation.
\end{proposition}

\section{Coordinate formulas}

With
\[
I=\begin{pmatrix}E&F\\F&G\end{pmatrix},
\qquad
II=\begin{pmatrix}e&f\\f&g\end{pmatrix},
\]
and \(S=I^{-1}II\), one obtains:

\begin{theorem}[Curvature formulas]\label{thm:vii17-formulas}
\[
K=\frac{eg-f^2}{EG-F^2},
\]
and
\[
H=\frac{Eg-2Ff+Ge}{2(EG-F^2)}.
\]
\end{theorem}

\begin{proof}
The first follows from \(\det(I^{-1}II)=\det II/\det I\). The second is one half of the trace of the \(2\times2\) matrix \(I^{-1}II\).
\end{proof}

\section{Point classification}

\begin{definition}[Elliptic, hyperbolic, parabolic, planar]\label{def:vii17-point-types}
At a point \(p\):
\begin{itemize}
\item \(K>0\): elliptic;
\item \(K<0\): hyperbolic;
\item \(K=0\) with \(S\neq0\): parabolic;
\item \(S=0\): planar.
\end{itemize}
\end{definition}

\begin{proposition}[Signs of normal curvature]\label{prop:vii17-signs}
At an elliptic point all nonzero normal curvatures have the same sign after choosing an orientation. At a hyperbolic point there are tangent directions with opposite signs of normal curvature.
\end{proposition}

\section{Asymptotic directions}

\begin{definition}[Asymptotic direction]\label{def:vii17-asymptotic}
A nonzero tangent direction \(v\) is asymptotic if
\[
II(v,v)=0,
\]
equivalently \(k_n(v)=0\).
\end{definition}

\begin{proposition}[Existence]\label{prop:vii17-asymptotic-existence}
At a hyperbolic point there are exactly two asymptotic lines in the tangent plane. At an elliptic point there are none.
\end{proposition}

\section{Umbilic points}

\begin{definition}[Umbilic]\label{def:vii17-umbilic}
A point is umbilic if
\[
S_p=\lambda\,\operatorname{id},
\]
equivalently \(k_1=k_2\).
\end{definition}

\begin{proposition}[Umbilic curvature relation]\label{prop:vii17-umbilic-relation}
At an umbilic point,
\[
K=H^2.
\]
In general,
\[
H^2-K=\frac{(k_1-k_2)^2}{4}\ge0.
\]
\end{proposition}

\section{Examples}

\begin{example}[Plane]\label{ex:vii17-plane}
A plane has \(k_1=k_2=0\), hence \(K=H=0\).
\end{example}

\begin{example}[Sphere]\label{ex:vii17-sphere}
For the outward normal under the convention \(S=-dN\),
\[
k_1=k_2=-\frac1R,\qquad K=\frac1{R^2},\qquad H=-\frac1R.
\]
Every point is umbilic.
\end{example}

\begin{example}[Cylinder]\label{ex:vii17-cylinder}
A circular cylinder has one zero principal curvature and one nonzero principal curvature. Hence \(K=0\) and \(H\neq0\).
\end{example}

\begin{example}[Saddle model]\label{ex:vii17-saddle}
Near a saddle point such as the origin of \(z=x^2-y^2\), the principal curvatures have opposite signs, so \(K<0\).
\end{example}

\section{Gauss-map interpretation}

\begin{theorem}[Jacobian of the Gauss map]\label{thm:vii17-gauss-jacobian}
Up to the orientation convention for area forms, the Jacobian determinant of the Gauss map is
\[
\det dN_p=\det(-S_p)=K(p).
\]
Thus Gaussian curvature measures infinitesimal area distortion of the Gauss map.
\end{theorem}

\section{What is intrinsic and what is extrinsic?}

The definition above uses the embedding in \(\mathbb R^3\), but Gaussian curvature is much more rigid than mean curvature. Theorema Egregium, developed later in the Riemannian arc, shows that \(K\) is determined by the first fundamental form alone. Mean curvature remains extrinsic.

\section{Exercises}
\begin{exercise}\label{exr:vii17-01}
State the principal curvatures in terms of the shape operator.
\end{exercise}
\begin{hint}
Use eigenvalues.
\end{hint}
\begin{solution}
They are precisely the two eigenvalues of the self-adjoint map \(S_p\).
\end{solution}

\begin{exercise}\label{exr:vii17-02}
Prove Euler's formula for normal curvature.
\end{exercise}
\begin{hint}
Use an orthonormal eigenbasis.
\end{hint}
\begin{solution}
Writing \(v=\cos\theta e_1+\sin\theta e_2\) and pairing \(Sv\) with \(v\) gives \(k_1\cos^2\theta+k_2\sin^2\theta\).
\end{solution}

\begin{exercise}\label{exr:vii17-03}
Why are principal directions orthogonal when \(k_1
e k_2\)?
\end{exercise}
\begin{hint}
Use self-adjointness.
\end{hint}
\begin{solution}
Eigenvectors of a self-adjoint map corresponding to distinct eigenvalues are orthogonal.
\end{solution}

\begin{exercise}\label{exr:vii17-04}
Show \(K=\det S\).
\end{exercise}
\begin{hint}
Use the eigenvalues.
\end{hint}
\begin{solution}
The determinant of a linear map equals the product of its eigenvalues, hence \(K=k_1k_2\).
\end{solution}

\begin{exercise}\label{exr:vii17-05}
Show \(H=\frac12\operatorname{tr}S\).
\end{exercise}
\begin{hint}
Use the eigenvalues.
\end{hint}
\begin{solution}
The trace is \(k_1+k_2\), so half the trace is the mean curvature.
\end{solution}

\begin{exercise}\label{exr:vii17-06}
Track \(K\) under orientation reversal.
\end{exercise}
\begin{hint}
Both eigenvalues change sign.
\end{hint}
\begin{solution}
Their product is unchanged, so \(K\) is orientation independent.
\end{solution}

\begin{exercise}\label{exr:vii17-07}
Track \(H\) under orientation reversal.
\end{exercise}
\begin{hint}
Both eigenvalues change sign.
\end{hint}
\begin{solution}
Their sum changes sign, so \(H\mapsto-H\).
\end{solution}

\begin{exercise}\label{exr:vii17-08}
Derive \(K=(eg-f^2)/(EG-F^2)\).
\end{exercise}
\begin{hint}
Use \(\det(I^{-1}II)\).
\end{hint}
\begin{solution}
One has \(K=\det II/\det I=(eg-f^2)/(EG-F^2)\).
\end{solution}

\begin{exercise}\label{exr:vii17-09}
Derive the coordinate formula for \(H\).
\end{exercise}
\begin{hint}
Compute the trace of \(I^{-1}II\).
\end{hint}
\begin{solution}
The trace equals \((Eg-2Ff+Ge)/(EG-F^2)\), and division by \(2\) gives \(H\).
\end{solution}

\begin{exercise}\label{exr:vii17-10}
Classify a point with \(K<0\).
\end{exercise}
\begin{hint}
Use the sign classification.
\end{hint}
\begin{solution}
It is hyperbolic.
\end{solution}

\begin{exercise}\label{exr:vii17-11}
Classify a point with \(K>0\).
\end{exercise}
\begin{hint}
Use the sign classification.
\end{hint}
\begin{solution}
It is elliptic.
\end{solution}

\begin{exercise}\label{exr:vii17-12}
Classify a point with \(K=0\) but \(S
eq0\).
\end{exercise}
\begin{hint}
One principal curvature vanishes.
\end{hint}
\begin{solution}
It is parabolic.
\end{solution}

\begin{exercise}\label{exr:vii17-13}
Explain why a plane is planar in the curvature classification.
\end{exercise}
\begin{hint}
Its shape operator vanishes.
\end{hint}
\begin{solution}
The normal is constant, so \(S=0\).
\end{solution}

\begin{exercise}\label{exr:vii17-14}
Show a sphere is umbilic.
\end{exercise}
\begin{hint}
Its shape operator is a scalar multiple of the identity.
\end{hint}
\begin{solution}
On a radius-\(R\) sphere, \(S=-(1/R)I\) for outward orientation.
\end{solution}

\begin{exercise}\label{exr:vii17-15}
Prove \(H^2-K=(k_1-k_2)^2/4\).
\end{exercise}
\begin{hint}
Expand both sides.
\end{hint}
\begin{solution}
From \(H=(k_1+k_2)/2\) and \(K=k_1k_2\), subtraction yields the stated square.
\end{solution}

\begin{exercise}\label{exr:vii17-16}
Deduce \(H^2\ge K\).
\end{exercise}
\begin{hint}
Use the previous identity.
\end{hint}
\begin{solution}
The right side is a square and is nonnegative.
\end{solution}

\begin{exercise}\label{exr:vii17-17}
When does equality \(H^2=K\) hold?
\end{exercise}
\begin{hint}
Set the square to zero.
\end{hint}
\begin{solution}
Exactly when \(k_1=k_2\), i.e. at umbilic points.
\end{solution}

\begin{exercise}\label{exr:vii17-18}
Define an asymptotic direction.
\end{exercise}
\begin{hint}
Use normal curvature.
\end{hint}
\begin{solution}
It is a nonzero tangent direction \(v\) with \(II(v,v)=0\), equivalently \(k_n(v)=0\).
\end{solution}

\begin{exercise}\label{exr:vii17-19}
Why do hyperbolic points have asymptotic directions?
\end{exercise}
\begin{hint}
Principal curvatures have opposite signs.
\end{hint}
\begin{solution}
Euler's formula varies continuously from one sign to the other and vanishes in two distinct lines.
\end{solution}

\begin{exercise}\label{exr:vii17-20}
Why do elliptic points have no asymptotic directions?
\end{exercise}
\begin{hint}
Principal curvatures have the same nonzero sign.
\end{hint}
\begin{solution}
Euler's formula is then a positive or negative weighted sum and cannot vanish.
\end{solution}

\begin{exercise}\label{exr:vii17-21}
Compute \(K\) for a cylinder.
\end{exercise}
\begin{hint}
One principal curvature is zero.
\end{hint}
\begin{solution}
Their product is zero.
\end{solution}

\begin{exercise}\label{exr:vii17-22}
Compute \(K,H\) for a plane.
\end{exercise}
\begin{hint}
Both principal curvatures vanish.
\end{hint}
\begin{solution}
Both invariants are zero.
\end{solution}

\begin{exercise}\label{exr:vii17-23}
Interpret \(K\) through the Gauss map.
\end{exercise}
\begin{hint}
Use \(\det dN\).
\end{hint}
\begin{solution}
It is the infinitesimal signed area distortion factor of the Gauss map, up to the standard orientation convention.
\end{solution}

\begin{exercise}\label{exr:vii17-24}
State the bridge to ruled/developable surfaces.
\end{exercise}
\begin{hint}
Use \(K=0\).
\end{hint}
\begin{solution}
VII/18 studies ruled surfaces and characterizes developable behavior through vanishing Gaussian curvature and degeneracy of the Gauss map.
\end{solution}

\section{Solved problem dossiers}
\begin{problem}[Legacy principal-curvature problem]\label{prob:vii17-01}
Show that principal curvatures are the extremal normal curvatures.
\end{problem}
\begin{solution}
Euler's formula expresses normal curvature as a convex combination of \(k_1,k_2\), so its extrema are the eigenvalues.
\end{solution}

\begin{problem}[Legacy Gaussian-curvature problem]\label{prob:vii17-02}
Derive the determinant formula for \(K\) in coordinates.
\end{problem}
\begin{solution}
Use \(S=I^{-1}II\), then take determinants.
\end{solution}

\begin{problem}[Legacy mean-curvature problem]\label{prob:vii17-03}
Derive the trace formula for \(H\) in coordinates.
\end{problem}
\begin{solution}
Invert the \(2\times2\) first-form matrix, multiply by \(II\), and take half the trace.
\end{solution}

\begin{problem}[Orientation-invariance problem]\label{prob:vii17-04}
Determine which standard curvature data survive orientation reversal.
\end{problem}
\begin{solution}
Principal curvatures and \(H\) reverse sign; \(K\), the unordered principal directions, and the point type determined by the sign of \(K\) do not.
\end{solution}

\begin{problem}[Umbilic problem]\label{prob:vii17-05}
Characterize umbilic points using the equality \(H^2=K\).
\end{problem}
\begin{solution}
The identity \(H^2-K=(k_1-k_2)^2/4\) shows equality exactly when \(k_1=k_2\).
\end{solution}

\begin{problem}[Hyperbolic-point problem]\label{prob:vii17-06}
Derive the existence of two asymptotic directions at a hyperbolic point.
\end{problem}
\begin{solution}
In a principal basis solve \(k_1\cos^2\theta+k_2\sin^2\theta=0\); opposite signs yield two tangent lines.
\end{solution}

\begin{problem}[Sphere problem]\label{prob:vii17-07}
Compute all principal, Gaussian, and mean curvature data for a sphere.
\end{problem}
\begin{solution}
Both principal curvatures equal \(-1/R\) for outward orientation, giving \(K=1/R^2\) and \(H=-1/R\).
\end{solution}

\begin{problem}[Cylinder problem]\label{prob:vii17-08}
Compute the curvature type of a circular cylinder.
\end{problem}
\begin{solution}
One principal curvature is zero and the other nonzero, so \(K=0\), \(H
e0\), and the points are parabolic.
\end{solution}

\begin{problem}[Plane problem]\label{prob:vii17-09}
Explain why every plane point is umbilic and planar.
\end{problem}
\begin{solution}
The zero operator is \(0\cdot I\), so the principal curvatures coincide at zero and \(S=0\).
\end{solution}

\begin{problem}[Gauss-map problem]\label{prob:vii17-10}
Relate \(K\) to the differential of the Gauss map.
\end{problem}
\begin{solution}
Because \(dN=-S\) in dimension two, \(\det dN=\det S=K\).
\end{solution}

\begin{problem}[Rayleigh-quotient problem]\label{prob:vii17-11}
Recover \(k_1,k_2\) from the normal-curvature function on the unit circle.
\end{problem}
\begin{solution}
They are its minimum and maximum; their eigendirections are the directions at which these extrema occur.
\end{solution}

\begin{problem}[Intrinsic-versus-extrinsic problem]\label{prob:vii17-12}
Explain the conceptual difference between \(K\) and \(H\).
\end{problem}
\begin{solution}
Both arise from the embedding through \(S\), but later Gauss theory proves \(K\) is determined by the first fundamental form, whereas \(H\) changes under isometric bendings in general.
\end{solution}

\section{Challenge problems}
\begin{problem}[Asymptotic-angle formula]\label{prob:vii17-ch01}
At a hyperbolic point derive the angle of an asymptotic direction relative to a principal direction.
\end{problem}
\begin{solution}
Solve \(k_1\cos^2\theta+k_2\sin^2\theta=0\). Thus \(\tan^2\theta=-k_1/k_2\), which is positive exactly when \(k_1k_2<0\).
\end{solution}

\begin{problem}[Recover principal curvatures from \(H,K\)]\label{prob:vii17-ch02}
Show that \(k_1,k_2\) are roots of a quadratic determined by \(H,K\).
\end{problem}
\begin{solution}
The characteristic polynomial is \(\lambda^2-2H\lambda+K\). Hence \(k_{1,2}=H\pm\sqrt{H^2-K}\).
\end{solution}

\begin{problem}[Gauss-map rank classification]\label{prob:vii17-ch03}
Relate the rank of \(dN_p\) to \(K\) and the point type.
\end{problem}
\begin{solution}
Since \(dN=-S\), rank two means \(K\ne0\); rank one means \(K=0\) but \(S\ne0\), a parabolic point; rank zero means \(S=0\), a planar point.
\end{solution}

\begin{problem}[Normal-curvature reconstruction]\label{prob:vii17-ch04}
Suppose the function \(v\mapsto k_n(v)\) is known for every unit tangent vector. Show it determines \(S_p\).
\end{problem}
\begin{solution}
The quadratic form \(q(v)=\langle Sv,v\rangle\) is known on the unit circle and hence by homogeneity on all vectors. Polarization recovers \(\langle Su,v\rangle\), and the metric then determines \(Su\) for every \(u\).
\end{solution}

\begin{problem}[Curvature discriminant]\label{prob:vii17-ch05}
Interpret \(H^2-K\) geometrically.
\end{problem}
\begin{solution}
It equals one quarter of the squared difference between principal curvatures. It measures anisotropy of normal bending and vanishes exactly at umbilics.
\end{solution}

\section*{Bridge to VII/18}
VII/18 applies the curvature machinery to ruled surfaces and isolates the developable case, where \(K=0\) and the Gauss map degenerates along rulings.
